%% ma: L10-B13-0096
%% lop: 10
%% chuong: 5
%% bai: 13
%% dang: 10.13.12
%% loai: TN4
%% muc: NB
%% dap_an: "B"
%% nguon: "chuyen-de-cac-so-dac-trung-cua-mau-so-lieu-khong-ghep-nhom-toan-10-chuong-trinh-moi.pdf (GV. Phan Nhật Linh) (P03), Bài 02 Dạng 2 Phần I Câu 24"
%% kien_thuc: "Lớp 10 Bài 13 (mốt)"
%% trang_thai: chinh-thuc
%% ghi_chu: "Tự duyệt theo yêu cầu giáo viên 22/10/2026."
\begin{ex}
Thời gian chạy $50$ m của $20$ học sinh được ghi lại trong bảng sau:
\begin{center}
\begin{tabular}{|c|c|c|c|c|c|}\hline
Thời gian (giây) & $8{,}3$ & $8{,}4$ & $8{,}5$ & $8{,}7$ & $8{,}8$\\\hline
Tần số & $2$ & $3$ & $9$ & $5$ & $1$\\\hline
\end{tabular}
\end{center}
Mốt của mẫu số liệu là
\choice{$M_o=9$}{\True $M_o=8{,}5$}{$M_o=8{,}7$}{$M_o=8{,}4$}
\loigiai{Giá trị $8{,}5$ có tần số lớn nhất ($9$) nên $M_o=8{,}5$. Chọn B.}
\end{ex}
